\begin{proof}[Proof of \cref{obj:thm:point-frontier}]
\begin{enumerate}
\item Choose the universal positive squared-risk constant from \cref{obj:lem:prior-reduction}, denoting it by \(c_{\mathrm{point}}\), and the universal upper-bound constant \(C_U>0\) from \cref{obj:lem:upper-risk}. Define
\[
c=\min\{c_{\mathrm{point}},C_U/2\},
\qquad C=C_U.
\]
Both entries in the minimum are positive, and therefore
\[
0<c\le C_U/2<C_U=C<\infty.
\]
These constants are independent of \(n,d,q\).
% lean: CausalSmith.Stat.MarNearcompleteFrontier.point_frontier

\item Fix positive integers \(n,d\) and \(q\in[1/2,1]\). Apply \cref{obj:lem:prior-reduction} with \(\alpha=1/4\in(0,1/2)\). Its combined squared-risk guarantee supplies
\[
c_{\mathrm{point}}\,r_{n,d,q}
\le \mathfrak R^*_{n,d,q}.
\]
By \cref{obj:def:rate},
\[
r_{n,d,q}=\frac1n+g_{n,d,q}^{\,2}\ge0.
\]
Consequently, the choice of \(c\) gives
\[
c\,r_{n,d,q}
\le c_{\mathrm{point}}\,r_{n,d,q}
\le \mathfrak R^*_{n,d,q}.
\]
% lean: CausalSmith.Stat.MarNearcompleteFrontier.prior_reduction

\item We verify that the estimator in \cref{obj:def:mm-estimator} belongs to \(\mathcal T_n\). Write
\[
o=(O_1,\ldots,O_n)\in O^n
\]
for the observed sample tuple; the range calculations below apply to every such tuple. The projection used there satisfies, for every real argument,
\[
-1\le \Pi_{[-1,1]}(t)
=\max\{-1,\min\{1,t\}\}\le1
\qquad(t\in\mathbb R).
\]
Thus the branch \(q=1\) takes values in \([-1,1]\), as does the fallback branch when \(q\ne1\) and either \(\ell_n<128\) or \(d\ge n\ell_n\).

For the remaining branch, define the Poisson weights and the averages over stream assignments by
\[
\omega_{n,N}
=e^{-n/2}\frac{(n/2)^N}{N!}
\qquad(N\in\{0,1,2,\ldots\}),
\]
\[
\overline\tau_N(o)
=4^{-N}
\sum_{\mathbf Z\in\{1,2,3,4\}^{N}}
\widetilde\tau(o;N,\mathbf Z)
\qquad(0\le N\le n).
\]
Here \(\widetilde\tau(o;N,\mathbf Z)\) is the projected randomized estimate specified in \cref{obj:def:mm-estimator} for the fixed sample, auxiliary count, and stream assignments. Each summand lies in \([-1,1]\), and there are \(4^N\) assignments. Hence
\[
-4^N
\le
\sum_{\mathbf Z\in\{1,2,3,4\}^{N}}
\widetilde\tau(o;N,\mathbf Z)
\le4^N,
\qquad
-1\le\overline\tau_N(o)\le1.
\]
For \(N=0\), the assignment set consists of the single empty tuple, so the same calculation applies.

The Poisson weights satisfy
\[
\omega_{n,N}\ge0,
\qquad
0\le\sum_{N=0}^{n}\omega_{n,N}
\le\sum_{N=0}^{\infty}\omega_{n,N}=1.
\]
Since the randomized estimate is zero when \(N>n\), its conditional average is
\[
\widehat\tau_{\mathrm{MM}}(o)
=\sum_{N=0}^{n}\omega_{n,N}\overline\tau_N(o)
\]
in this branch. Multiplying the bounds on each average by its nonnegative weight and summing gives
\[
-1
\le-\sum_{N=0}^{n}\omega_{n,N}
\le\widehat\tau_{\mathrm{MM}}(o)
\le\sum_{N=0}^{n}\omega_{n,N}
\le1.
\]
The projected formulas and finite averages are measurable. Thus, on the entire Cartesian sample space,
\[
\widehat\tau_{\mathrm{MM}}\in\mathcal T_n,
\]
with \(\mathcal T_n\) as defined in \cref{obj:synth_6}.
% lean: CausalSmith.Stat.MarNearcompleteFrontier.tauhatMM_range

\item The zero-pair-mass specialization of the construction in \cref{obj:def:paired-prior-handle} places all normalized mass at its filler label. It has independent fair treatment and complete arrival, and supplies
\[
\mathcal M(d,q)\ne\varnothing
\qquad(d\ge1,\ q\in[1/2,1]),
\]
including \(d=1\). For the real suprema below, we use the conditional real supremum: a nonempty set bounded above has its usual least upper bound, while a set unbounded above is assigned supremum zero. Fix \(T\in\mathcal T_n\). If its set of squared-error risks is bounded above, nonnegativity of squared loss and the supremum property give
\[
0\le
\mathbb E_P\!\left[
\{T(o)-\tau(P)\}^{2}
\right]
\le
\sup_{P'\in\mathcal M(d,q)}
\mathbb E_{P'}\!\left[
\{T(o)-\tau(P')\}^{2}
\right]
\qquad(P\in\mathcal M(d,q)).
\]
The law class is nonempty, so this bounds the supremum below by zero. If the set of squared-error risks is unbounded above, the stated convention gives
\[
\sup_{P\in\mathcal M(d,q)}
\mathbb E_P\!\left[
\{T(o)-\tau(P)\}^{2}
\right]=0.
\]
In either case,
\[
0\le
\sup_{P\in\mathcal M(d,q)}
\mathbb E_P\!\left[
\{T(o)-\tau(P)\}^{2}
\right]
\qquad(T\in\mathcal T_n).
\]
The worst-case risks are therefore bounded below by zero. Evaluating the infimum in \cref{obj:def:point-risk} at the admissible estimator established above yields
\[
\mathfrak R^*_{n,d,q}
\le
\sup_{P\in\mathcal M(d,q)}
\mathbb E_P\!\left[
\{\widehat\tau_{\mathrm{MM}}(o)-\tau(P)\}^{2}
\right].
\]
For every law in this supremum, the sample consists of independent records with its observed-law distribution. Consequently, \cref{obj:lem:upper-risk} gives
\[
\sup_{P\in\mathcal M(d,q)}
\mathbb E_P\!\left[
\{\widehat\tau_{\mathrm{MM}}(o)-\tau(P)\}^{2}
\right]
\le C_U r_{n,d,q}
=C r_{n,d,q}.
\]
Together with the lower bound, this proves the claimed inequalities throughout the stated parameter range.
% lean: CausalSmith.Stat.MarNearcompleteFrontier.point_frontier
\end{enumerate}
\end{proof}
